\documentclass[12pt, a4paper]{article}

% ── PACKAGES ──────────────────────────────────────────────────
\usepackage[utf8]{inputenc}
\usepackage[T1]{fontenc}
\usepackage{geometry}
\geometry{margin=2.5cm}
\usepackage{hyperref}
\hypersetup{
    colorlinks=true,
    linkcolor=blue,
    urlcolor=blue,
    citecolor=blue,
    pdftitle={Androgen Receptor as a Continuous Epigenetic
              Depth Axis in Triple-Negative Breast Cancer:
              A Geometry-Derived Reframing of a Binary
              Classifier},
    pdfauthor={Eric Robert Lawson},
    pdfkeywords={androgen receptor, AR, TNBC,
                 triple-negative breast cancer,
                 continuous depth axis, EZH2, epigenetic
                 lock, Waddington landscape, attractor
                 geometry, LAR, luminal androgen receptor,
                 GSE176078, single cell, Pearson correlation,
                 OrganismCore, biomarker, depth score,
                 G-1 enobosarm, AR targeting, patient
                 selection}
}
\usepackage{amsmath}
\usepackage{booktabs}
\usepackage{longtable}
\usepackage{array}
\usepackage{parskip}
\usepackage{microtype}
\usepackage{authblk}
\usepackage{titlesec}
\usepackage{fancyhdr}
\usepackage{xcolor}
\usepackage{float}
\usepackage{enumitem}
\usepackage{setspace}

% ── COLOURS ───────────────────────────────────────────────────
\definecolor{repobox}{RGB}{240,247,255}
\definecolor{findingbox}{RGB}{245,255,245}
\definecolor{novelbox}{RGB}{255,248,220}
\definecolor{mechanismbox}{RGB}{250,245,255}
\definecolor{actionbox}{RGB}{230,255,230}
\definecolor{contrastbox}{RGB}{248,240,255}
\definecolor{warnbox}{RGB}{255,235,235}

% ── HEADER / FOOTER ───────────────────────────────────────────
\pagestyle{fancy}
\fancyhf{}
\lhead{\footnotesize CS-LIT-6: AR as Continuous Depth\\
       \footnotesize Axis in TNBC}
\rhead{\footnotesize OrganismCore \textbar{} 2026-03-06}
\rfoot{\small Page \thepage}
\renewcommand{\headrulewidth}{0.4pt}
\setlength{\headheight}{24pt}
\addtolength{\topmargin}{-10pt}

% ── SECTION FORMATTING ────────────────────────────────────────
\titleformat{\section}
  {\large\bfseries}{\thesection}{1em}{}[\titlerule]
\titleformat{\subsection}
  {\normalsize\bfseries}{\thesubsection}{1em}{}
\titleformat{\subsubsection}
  {\small\bfseries}{\thesubsubsection}{1em}{}

% ══════════════════════════════════════════════════════════════
% TITLE
% ══════════════════════════════════════════════════════════════
\title{
    \vspace{-1em}
    \textbf{Androgen Receptor as a Continuous
    Epigenetic Depth Axis in Triple-Negative
    Breast Cancer: A Geometry-Derived Reframing
    of a Binary Classifier}\\[0.5em]
    \large AR Expression Is Not a Subtype
    Identity Marker Within TNBC --- It Is a
    Quantitative Measure of How Deeply the
    Tumor Is Locked in the EZH2 Attractor\\[0.4em]
    \normalsize \textit{OrganismCore Framework ---
    Document CS-LIT-6}\\[0.3em]
    \small Pearson $r = -0.378$,
    $p = 6.23 \times 10^{-147}$,
    $n = 4{,}312$ TNBC cells (GSE176078)\\[0.3em]
    \small Part of the BRCA Cross-Subtype Analysis
    (BRCA-S8h, 2026-03-05)
}

\author[1]{Eric Robert Lawson}
\affil[1]{
    OrganismCore\\
    \href{mailto:OrganismCore@proton.me}
         {OrganismCore@proton.me}\\
    ORCID:
    \href{https://orcid.org/0009-0002-0414-6544}
         {0009-0002-0414-6544}
}

\date{March 6, 2026}

% ══════════════════════════════════════════════════════════════
\begin{document}
% ══════════════════════════════════════════════════════════════

\maketitle
\thispagestyle{fancy}

% ── REPOSITORY POINTER ────────────────────────────────────────
\begin{center}
\colorbox{repobox}{
\begin{minipage}{0.88\textwidth}
\vspace{0.5em}
\centering
\textbf{Full computational record --- scripts, data
caches, reasoning artifacts, and raw logs:}\\[0.4em]
\href{https://github.com/Eric-Robert-Lawson/attractor-oncology}
{\texttt{https://github.com/Eric-Robert-Lawson/attractor-oncology}}
\\[0.3em]
\small Part of the OrganismCore BRCA Cross-Subtype
Analysis (BRCA-S8h, 2026-03-05).\\
All predictions were locked from attractor geometry
\textbf{before} any literature review was performed.
\vspace{0.5em}
\end{minipage}}
\end{center}

\vspace{1em}

% ── FINDING BOX ───────────────────────────────────────────────
\begin{center}
\colorbox{findingbox}{
\begin{minipage}{0.88\textwidth}
\vspace{0.5em}
\textbf{THE FINDING, STATED PRECISELY:}\\[0.4em]
Within the TNBC cell population of GSE176078
($n = 4{,}312$ TNBC cells), androgen receptor
(AR) expression forms a strong continuous
inverse correlation with the EZH2-driven
attractor lock depth:

\[
r = -0.378,\quad
p = 6.23 \times 10^{-147},\quad
n = 4{,}312
\]

AR is not acting as a subtype identity marker
within TNBC. It is acting as a continuous
quantitative readout of how far the tumor cell
has descended into the EZH2 epigenetic attractor.
Lower AR = deeper lock. Higher AR = shallower
lock. The relationship is graded, not binary.
\vspace{0.5em}
\end{minipage}}
\end{center}

\vspace{1em}

% ══════════════════════════════════════════════════════════════
\begin{abstract}
% ══════════════════════════════════════════════════════════════
Androgen receptor (AR) expression in triple-negative
breast cancer (TNBC) is currently classified in the
literature as a binary marker: AR-positive tumors
belong to the Luminal Androgen Receptor (LAR)
subtype; AR-negative tumors do not. This binary
framework treats a continuous biological variable as
categorical. We report a geometry-derived finding
from Waddington attractor analysis of 19,542 single
cancer cells (GSE176078) that fundamentally reframes
this interpretation. Within the 4,312 TNBC cells in
this dataset, AR expression forms a strong,
continuous, statistically significant inverse
correlation with EZH2-driven epigenetic lock depth
($r = -0.378$, $p = 6.23 \times 10^{-147}$,
$n = 4{,}312$). This is not a subtype classification
signal. AR level within the TNBC cell population
is a quantitative readout of attractor position:
lower AR = deeper EZH2 lock; higher AR = shallower
lock, with residual luminal identity program
activity. This reframing has three direct clinical
consequences. First, AR level continuously predicts
the depth of the epigenetic lock and therefore the
magnitude of benefit expected from EZH2-targeting
interventions. Second, it explains why AR-targeting
therapies (G-1/enobosarm) fail in deeply
de-differentiated TNBC: these cells have already
lost AR expression as a consequence of deep lock,
not as a targetable feature. Third, it establishes
AR as a component of the TNBC Depth Score rather
than as an independent binary classifier. The
finding is novel. The published literature does
not frame AR as a continuous epigenetic depth
axis within TNBC, and the specific correlation
statistics reported here have no published
equivalent.
\end{abstract}

\tableofcontents
\newpage

% ══════════════════════════════════════════════════════════════
\section{Context and Origin}
% ══════════════════════════════════════════════════════════════

This document covers CS-LIT-6 from the OrganismCore
BRCA Cross-Subtype Literature Check (BRCA-S8h,
2026-03-05):

\begin{quote}
\textbf{CS-LIT-6}: AR as a continuous depth axis
within TNBC ($r = -0.378$,
$p = 6.23 \times 10^{-147}$, $n = 4{,}312$ cells)
\end{quote}

The finding was derived from attractor geometry
before any literature review was performed. The
literature check assigned it a \textbf{NOVEL}
verdict: the specific framing and statistics
have no direct published equivalent.

\subsection{Relationship to Other Published
Documents in This Series}

\begin{itemize}
    \item \textbf{CS-LIT-2 (Six Lock Type
    Classification)}\\
    \href{https://doi.org/10.5281/zenodo.18884158}
    {DOI: 10.5281/zenodo.18884158}\\
    TNBC is Lock Type 3 (deep EZH2/PRC2 lock).
    CS-LIT-6 provides the continuous within-TNBC
    geometry that underlies Lock Type 3.

    \item \textbf{CS-LIT-1 (FOXA1/EZH2 Ratio)}\\
    \href{https://doi.org/10.5281/zenodo.18883922}
    {DOI: 10.5281/zenodo.18883922}\\
    The FOXA1/EZH2 ratio orders subtypes
    \textit{across} TNBC. CS-LIT-6 demonstrates
    that AR performs an analogous continuous
    ordering function \textit{within} TNBC.

    \item \textbf{CS-LIT-11/25 (TNBC Depth Score)}\\
    \href{https://doi.org/10.5281/zenodo.18891523}
    {DOI: 10.5281/zenodo.18891523}\\
    The TNBC Depth Score uses AR as one of three
    components. CS-LIT-6 is the foundational
    document establishing AR's role as a
    continuous component rather than a binary one.

    \item \textbf{CS-LIT-10/24 (EZH2 Paradox)}\\
    \href{https://doi.org/10.5281/zenodo.18891318}
    {DOI: 10.5281/zenodo.18891318}\\
    The EZH2 paradox describes what happens when
    the deep lock is disrupted by chemotherapy.
    CS-LIT-6 quantifies the depth of that lock
    using AR as the continuous readout.
\end{itemize}

% ══════════════════════════════════════════════════════════════
\section{The Problem with the Binary AR Framework}
% ══════════════════════════════════════════════════════════════

\subsection{What the Published Literature Says}

The published literature classifies TNBC using
AR as a binary marker:

\begin{itemize}
    \item \textbf{AR-positive TNBC}
    ($\geq 1\%$ or $\geq 10\%$ nuclear staining
    by IHC, depending on study): classified as
    the Luminal Androgen Receptor (LAR) subtype.
    Characterised by luminal gene expression,
    lower proliferation, lower pCR rates,
    and potential sensitivity to AR-targeting
    therapies.

    \item \textbf{AR-negative TNBC}: All
    remaining TNBC. Treated as a single category
    by the binary framework.
\end{itemize}

This binary framing drives clinical decisions:
AR-positive TNBC patients are candidates for
AR-targeting therapies (enzalutamide,
bicalutamide, G-1/enobosarm); AR-negative
patients are not.

\subsection{Why the Binary Framework Is Incomplete}

The binary AR framework implicitly treats a
continuous biological variable as a categorical
one. The question it cannot answer is:

\begin{center}
\colorbox{contrastbox}{
\begin{minipage}{0.82\textwidth}
\vspace{0.4em}
\textit{Within the AR-negative population, are
all tumors equivalent? Or does the level of AR
suppression --- how far below the detection
threshold AR expression falls --- carry
biological information?}
\vspace{0.4em}
\end{minipage}}
\end{center}

The attractor geometry analysis answers this
directly: AR suppression is graded. The deeper
the EZH2 lock, the lower the AR expression ---
continuously, not categorically. Tumors do not
simply turn AR off. They suppress it
progressively as lock depth increases.

% ══════════════════════════════════════════════════════════════
\section{The Geometric Finding}
% ══════════════════════════════════════════════════════════════

\subsection{Data and Method}

\begin{center}
\colorbox{repobox}{
\begin{minipage}{0.82\textwidth}
\vspace{0.4em}
\textbf{Dataset:} GSE176078\\
\textbf{Population:} 4,312 TNBC cells (subset
of 19,542 total single cancer cells)\\
\textbf{Method:} Waddington attractor geometry;
within-population Pearson correlation between
AR expression and the EZH2-driven attractor
depth axis\\
\textbf{Analysis type:} Within-population
continuous correlation (see CS-LIT-27:
within-population Pearson $r$ as a circuit
integrity test)
\vspace{0.4em}
\end{minipage}}
\end{center}

\vspace{0.5em}

\subsection{The Result}

\begin{center}
\colorbox{findingbox}{
\begin{minipage}{0.88\textwidth}
\vspace{0.6em}
\[
r = -0.378,\quad
p = 6.23 \times 10^{-147},\quad
n = 4{,}312
\]
\vspace{0.3em}
Within the TNBC cell population, AR expression
is strongly and significantly inversely
correlated with EZH2-driven attractor lock depth.
The relationship is continuous, graded, and
highly statistically significant.\\[0.4em]
\textbf{Direction:} Lower AR $\rightarrow$
deeper EZH2 lock.
Higher AR $\rightarrow$ shallower lock,
with residual luminal identity activity.
\vspace{0.5em}
\end{minipage}}
\end{center}

\subsection{\texorpdfstring{Why $p = 6.23 \times 10^{-147}$
Is Not Inflated}{Why the p-value Is Not Inflated}}

A $p$-value of $6.23 \times 10^{-147}$ reflects
a large single-cell dataset ($n = 4{,}312$) with
a moderate correlation ($r = -0.378$). The
extreme $p$-value is driven by sample size, not
by an artefactually large effect. The $r$ value
of $-0.378$ is a moderate correlation: meaningful
and consistent, but not claiming a deterministic
one-to-one relationship. Many other factors
modulate AR expression within individual cells.
The claim is that the relationship is real,
graded, and consistent --- not that AR is a
perfect proxy for lock depth.

\subsection{What the Geometry Shows}

The Waddington landscape of TNBC cells maps the
EZH2/PRC2 epigenetic axis as the primary
determinant of attractor position within TNBC.
Along this axis:

\begin{itemize}
    \item Cells at the deepest attractor position
    (maximum EZH2 activity, FOXA1 fully silenced,
    maximum de-differentiation) express near-zero
    AR.

    \item Cells at shallower attractor positions
    (moderate EZH2 activity, partial FOXA1
    suppression, partial luminal identity
    retention) express detectable AR.

    \item The transition between these states
    is continuous across the cell population ---
    not a binary switch.
\end{itemize}

The LAR subtype, as defined by the binary
framework, corresponds to the shallow end of
this continuous distribution --- cells that
have not descended deeply into the EZH2
attractor. They are not a categorically
different entity. They are the least-locked
end of the TNBC attractor spectrum.

% ══════════════════════════════════════════════════════════════
\section{The Three Clinical Consequences}
% ══════════════════════════════════════════════════════════════

\subsection{Consequence 1: AR Level Predicts
EZH2 Inhibitor Benefit Magnitude}

If AR is a continuous readout of lock depth,
then AR level predicts how much EZH2 inhibitor
benefit a patient will receive:

\begin{center}
\begin{tabular}{p{4cm}p{4cm}p{4.5cm}}
\toprule
\textbf{AR Level} &
\textbf{Attractor Position} &
\textbf{Expected EZH2i Benefit} \\
\midrule
Low / absent &
Deep lock. EZH2 maximally active. &
Highest benefit from tazemetostat.
Attractor is deepest; inhibition
has maximum effect. \\[0.5em]
Moderate &
Intermediate lock. &
Moderate benefit. \\[0.5em]
High (LAR) &
Shallow lock. Partial luminal
identity retained. &
Lower EZH2i benefit. PRC2 less
dominant. Alternative pathways
more relevant. \\
\bottomrule
\end{tabular}
\end{center}

\vspace{0.5em}

This is a dose--response prediction: the
\textit{magnitude} of EZH2 inhibitor benefit
is graded by AR level in inverse proportion.
The binary AR-positive/AR-negative framework
cannot capture this gradient.

\subsection{Consequence 2: Explanation for
G-1/Enobosarm Failure in AR-Low TNBC
(CS-LIT-12)}

AR-targeting therapies (G-1/enobosarm,
enzalutamide, bicalutamide) require AR
expression as their target. The geometric
finding provides the causal explanation for
why these therapies fail in deeply
de-differentiated TNBC:

\begin{center}
\colorbox{warnbox}{
\begin{minipage}{0.88\textwidth}
\vspace{0.5em}
\textbf{The G-1 failure mechanism:}\\[0.4em]
In deeply locked TNBC (AR-low), AR is not
absent because of a mutation or deletion.
It is absent because the EZH2/PRC2 complex
has silenced the luminal identity program
that would express it. AR is a downstream
casualty of lock depth, not an independent
targetable feature.\\[0.4em]
Administering an AR agonist (G-1) or antagonist
to cells with near-zero AR expression is
targeting an absent protein. The therapeutic
premise of AR-targeting in this population
is invalid --- not because AR-targeting does
not work, but because the target has already
been removed by the very epigenetic state
that characterises these tumors.\\[0.4em]
This explains the treatment-context dependence
of AR prognosis in TNBC observed in clinical
trials (CS-LIT-12 in BRCA-S8h): AR status
predicts response only when AR is present
as an expressed, functional protein. In
deeply locked TNBC, AR expression has already
been lost.
\vspace{0.5em}
\end{minipage}}
\end{center}

\subsection{Consequence 3: AR as a Component
of the TNBC Depth Score}

The continuous AR--lock depth relationship
directly justifies AR's inclusion as a
component of the TNBC Depth Score
(CS-LIT-11/25,
\href{https://doi.org/10.5281/zenodo.18891523}
{DOI: 10.5281/zenodo.18891523}).

In that instrument, AR is not used as a binary
classifier (AR-positive vs AR-negative). It is
used as a continuous variable: the lower the AR
expression, the greater its contribution to the
depth score. This is only valid if the continuous
AR--depth relationship demonstrated here is real.
CS-LIT-6 is therefore the foundational
justification for how AR is used in the
depth score.

% ══════════════════════════════════════════════════════════════
\section{What the Published Literature Does and
Does Not Contain}
% ══════════════════════════════════════════════════════════════

\subsection{What IS in the Published Literature}

\begin{itemize}
    \item AR expression defines the LAR subtype
    of TNBC (binary, $\geq$1\% or $\geq$10\%
    IHC).
    \item LAR TNBC has lower pCR rates and lower
    immune infiltration compared to non-LAR TNBC.
    \item AR-targeting therapies are being tested
    in LAR/AR-positive TNBC.
    \item EZH2 is overexpressed in TNBC and is
    associated with poor prognosis.
    \item AR and EZH2 are both recognised
    molecular features of TNBC.
\end{itemize}

\subsection{What is NOT in the Published
Literature}

\begin{itemize}
    \item The continuous within-population
    inverse correlation between AR expression
    and EZH2-driven attractor lock depth in
    TNBC cells.
    \item The framing of AR as a continuous
    epigenetic depth axis (rather than a binary
    subtype marker) within TNBC.
    \item The specific statistics
    ($r = -0.378$, $p = 6.23 \times 10^{-147}$,
    $n = 4{,}312$) reported here.
    \item The causal interpretation that AR
    suppression is a \textit{consequence} of
    EZH2 lock depth, not an independent
    categorical classifier.
    \item The derived prediction that EZH2
    inhibitor benefit magnitude is inversely
    graded by AR level within TNBC.
\end{itemize}

\begin{center}
\colorbox{novelbox}{
\begin{minipage}{0.88\textwidth}
\vspace{0.5em}
\textbf{VERDICT: NOVEL}\\[0.3em]
The continuous AR--EZH2 lock depth axis in
TNBC, framed geometrically using single-cell
Waddington attractor analysis, with the specific
correlation statistics reported here, has no
direct published equivalent as of 2026-03-06.
The binary LAR/non-LAR framework is well
established in the literature. The reframing
of that binary as a continuous attractor
depth gradient is not.
\vspace{0.5em}
\end{minipage}}
\end{center}

% ══════════════════════════════════════════════════════════════
\section{Comparison: FOXA1 Across Subtypes vs
AR Within TNBC}
% ══════════════════════════════════════════════════════════════

CS-LIT-5 (BRCA-S8h) documents that FOXA1 acts
as a continuous depth axis \textit{across}
breast cancer subtypes --- not as a binary
identity marker between subtypes.
CS-LIT-6 documents the analogous finding for AR
\textit{within} TNBC specifically.

The structural parallel is exact:

\begin{center}
\begin{tabular}{p{4.5cm}p{4.5cm}p{3cm}}
\toprule
\textbf{Feature} &
\textbf{Standard view} &
\textbf{Geometric reframing} \\
\midrule
FOXA1 (across subtypes) &
Binary: luminal vs non-luminal &
Continuous depth axis
across all six subtypes \\[0.5em]
AR (within TNBC) &
Binary: LAR vs non-LAR &
Continuous depth axis
within the TNBC attractor \\
\bottomrule
\end{tabular}
\end{center}

\vspace{0.5em}

Both findings arise from the same methodological
principle: when a binary classifier is examined
within a continuous Waddington landscape, the
binary threshold is revealed to be an arbitrary
cut on a graded biological spectrum.

% ══════════════════════════════════════════════════════════════
\section{Validation Pathway}
% ══════════════════════════════════��═══════════════════════════

\begin{center}
\colorbox{actionbox}{
\begin{minipage}{0.88\textwidth}
\vspace{0.5em}
\textbf{Proposed analyses:}\\[0.4em]
\textbf{Analysis 1 --- Independent single-cell
replication:}\\
Reproduce the within-TNBC AR--EZH2 correlation
in an independent single-cell TNBC dataset.
Confirm that the relationship is continuous,
not bimodal.\\[0.5em]
\textbf{Analysis 2 --- Bulk RNA-seq
confirmation:}\\
Confirm the inverse AR--depth relationship in
a bulk RNA TNBC cohort (e.g., TCGA-BRCA TNBC
subset or METABRIC). Compute the correlation
between AR expression and an EZH2-activity
proxy (e.g., EZH2 expression or H3K27me3
surrogate signature).\\[0.5em]
\textbf{Analysis 3 --- Clinical dose--response
test:}\\
In a TNBC cohort with AR IHC and tazemetostat
treatment data: test whether tazemetostat
benefit (response rate, PFS) is inversely
graded by pre-treatment AR H-score. This is
the falsifiable prediction from the continuous
depth axis model.\\[0.5em]
\textbf{Analysis 4 --- G-1 failure mechanism:}\\
In a TNBC cohort with AR IHC and G-1/enobosarm
or enzalutamide data: confirm that the lowest
AR expressors show no benefit, consistent with
CS-LIT-12.
\vspace{0.5em}
\end{minipage}}
\end{center}

% ══════════════════════════════════════════════════════════════
\section{What This Document Claims and Does
Not Claim}
% ══════════════════════════════════════════════════════════════

\subsection{What IS Claimed}

\begin{itemize}
    \item Within 4,312 TNBC cells in GSE176078,
    AR expression is continuously and inversely
    correlated with EZH2-driven attractor lock
    depth ($r = -0.378$,
    $p = 6.23 \times 10^{-147}$).
    \item The correlation is statistically
    robust and geometrically consistent with
    the Waddington landscape model.
    \item The binary LAR/non-LAR classification
    is a threshold imposed on a continuous
    biological variable.
    \item AR level predicts EZH2 inhibitor
    benefit magnitude in an inverse relationship.
    \item AR suppression in deeply locked TNBC
    is a consequence of lock depth, not an
    independent categorical feature.
    \item This framing has no direct published
    equivalent as of 2026-03-06.
\end{itemize}

\subsection{What is NOT Claimed}

\begin{itemize}
    \item The continuous AR axis has not been
    validated prospectively or in a clinical
    trial.
    \item The dose--response prediction for
    EZH2i benefit has not been tested.
    \item The LAR subtype designation is not
    invalidated. It remains clinically useful
    as a threshold. The claim is that treating
    it as a continuous axis captures additional
    information.
    \item The finding is from a single
    single-cell dataset and requires replication
    in independent data.
\end{itemize}

% ═════════════════════════════════════════���════════════════════
\section{Locked Statement}
% ══════════════════════════════════════════════════════════════

\begin{center}
\colorbox{findingbox}{
\begin{minipage}{0.88\textwidth}
\vspace{0.6em}
\textbf{Stated once, precisely, without inflation:}
\\[0.5em]
Within 4,312 TNBC cells in GSE176078,
AR expression is continuously and inversely
correlated with EZH2-driven epigenetic lock
depth ($r = -0.378$, $p = 6.23 \times 10^{-147}$,
$n = 4{,}312$). The relationship is graded, not
binary. The standard LAR/non-LAR binary
classification is a threshold on this continuous
spectrum. AR level within TNBC is a quantitative
readout of how deeply the tumor cell population
is locked in the EZH2/PRC2 attractor state.
Lower AR = deeper lock = higher TNBC Depth Score
= greater expected benefit from EZH2-targeting
intervention. Higher AR = shallower lock =
partial luminal identity retained = candidate
for AR-directed therapies. The causal direction
is: EZH2 lock depth suppresses AR, not the
reverse. This framing and these statistics have
no direct published equivalent as of 2026-03-06.
\vspace{0.5em}
\end{minipage}}
\end{center}

% ══════════════════════════════════════════════════════════════
\section{Document Metadata}
% ══════════════════════════════════════════════════════════════

\begin{tabular}{ll}
\toprule
\textbf{Field} & \textbf{Value} \\
\midrule
Document ID & CS-LIT-6 \\
Parent document & BRCA-S8h (2026-03-05) \\
Verdict & NOVEL \\
Date & 2026-03-06 \\
Author & Eric Robert Lawson \\
ORCID & 0009-0002-0414-6544 \\
Contact & OrganismCore@proton.me \\
Repository &
\href{https://github.com/Eric-Robert-Lawson/attractor-oncology}
{attractor-oncology} \\
Source data & GSE176078 ($n=4{,}312$ TNBC cells) \\
Key statistic &
$r=-0.378$, $p=6.23\times10^{-147}$,
$n=4{,}312$ \\
Six Lock Type DOI &
\href{https://doi.org/10.5281/zenodo.18884158}
{10.5281/zenodo.18884158} \\
FOXA1/EZH2 Ratio DOI &
\href{https://doi.org/10.5281/zenodo.18883922}
{10.5281/zenodo.18883922} \\
TNBC Depth Score DOI &
\href{https://doi.org/10.5281/zenodo.18891523}
{10.5281/zenodo.18891523} \\
EZH2 Paradox DOI &
\href{https://doi.org/10.5281/zenodo.18891318}
{10.5281/zenodo.18891318} \\
License & CC BY 4.0 \\
\bottomrule
\end{tabular}

\end{document}
